% PROVENANCE: independently audited by root 2026-11-19 14:09 UTC; incorporated into part_B.tex.
\paragraph*{Archived drafting provenance; authoritative version is in part B.}
\begin{saresult}{B.041}{Proved}{Every infinite valuation domain and every torsion-free extension}
Let $V$ be an infinite valuation domain with fraction field $K$,
and let $V\subset S$ be any $V$-torsion-free unital extension.
Then every polynomial tensor map for $\Omega(V,S)$ is an
isomorphism, and the unique compatible plethory exists.
More precisely, put $L=S_{\mathrm{rat}}$. If the residue field
of $V$ is infinite or its maximal ideal is nonprincipal, then
$\Omega(V,S)=V[X]$. If its nonzero maximal ideal is principal,
with finite residue field, then
\[
 \Omega(V,S)=
 \begin{cases}
 V[X],&L=V,\\
 \operatorname{Int}(V),&L\ne V,
 \end{cases}
\]
and $\operatorname{Int}(V)$ has an explicit free Newton basis.
This covers valuation domains of arbitrary rank and does not
assume Noetherianity. A finite valuation domain is a finite field
and is governed instead by B.013.
\end{saresult}
\begin{proof}
The field case is B.014, so suppose $V$ is not a field.
B.014 reduces $S$ to the overring $L\subset K$.
An infinite residue field gives $\operatorname{Int}(V)=V[X]$
by the invertible Vandermonde argument; hence it gives the
asserted equality for every coefficient overring.

Suppose the maximal ideal is nonprincipal. The ordered value
group $\Gamma$ has no least positive element: a least positive
value would give a generator of that ideal, and conversely.
For every $\gamma>0$ and every positive integer $N$, there is
$\delta>0$ with $N\delta<\gamma$. To verify this without an
Archimedean assumption, choose $0<\epsilon<\gamma$; any positive
$\epsilon$ has a positive $\eta$ with $2\eta\le\epsilon$.
Indeed choose $0<\xi<\epsilon$ and take $\eta=\xi$ if
$2\xi\le\epsilon$, otherwise take $\eta=\epsilon-\xi$.
Iterate until $2^j\ge N$.

Take $F\in\operatorname{Int}(V)$ of degree $d\ge1$. If one
coefficient had value $-\gamma<0$, choose $a\in V$ with value
$\delta>0$ so small that
\[
 \frac{d(d+1)(d+2)}6\,\delta<\gamma.
\]
The $d+1$ points $0,a,a^2,\ldots,a^d$ are distinct. Their
Vandermonde determinant has value exactly
$d(d+1)(d+2)\delta/6$: the difference of distinct positive
powers of $a$ has the smaller of their two values.
Cramer's rule, with all values $F(a^i)$ in $V$, forces every
coefficient to have value at least the negative of this
quantity, contradicting $-\gamma$. Constant polynomials are
already in $V$. Thus again $\operatorname{Int}(V)=V[X]$.

It remains to suppose the maximal ideal is $\pi V$ and the
residue field has $q$ elements. Choose representatives
$c_0=0,c_1,\ldots,c_{q-1}\in V$. For the base-$q$ expansion
$n=\sum_j n_jq^j$ put $a_n=\sum_jc_{n_j}\pi^j$ and
\[
 w_q(n)=\sum_{j\ge1}\lfloor n/q^j\rfloor,\qquad
 b_n(X)=\frac{\prod_{i<n}(X-a_i)}{\prod_{i<n}(a_n-a_i)}.
\]
The denominator has value exactly $w_q(n)v(\pi)$, by the
first differing digit of each pair. Every $x\in V$ has any
prescribed finite number of successive residue digits, since
its maximal ideal is $\pi V$. Among $a_0,\ldots,a_{n-1}$ at
least $\lfloor n/q^j\rfloor$ share the first $j$ digits of
$x$. Summing these finite lower bounds shows that the
numerator has value at least $w_q(n)v(\pi)$.
This argument uses only finitely many digit levels with
nonzero lower bounds, so it remains valid in a higher-rank
value group. Hence $b_n\in\operatorname{Int}(V)$ and
$b_n(a_n)=1$. Successive evaluation at these nodes gives
\[
 \operatorname{Int}(V)=\bigoplus_{n\ge0}Vb_n.
\]
Every $b_n$ has coefficients in $V[\pi^{-1}]$.

If $L\ne V$, choose $b\in L\setminus V$.
The least positive value $v(\pi)$ satisfies
$-v(b)\ge v(\pi)$. Therefore $(\pi b)^{-1}\in V$ and
$\pi^{-1}=(\pi b)^{-1}b\in L$.
The entire Newton basis belongs to $L[X]$, so
$\Omega(V,L)=\operatorname{Int}(V)$. If $L=V$ the module is
$V[X]$. These modules are free over $V$ in every case;
B.018 gives all tensor isomorphisms and the plethory.
\end{proof}
